% !TeX root = ../main.tex
\section{Introduction}
\label{sec:intro}

% Large language model (LLM) multi-agent systems coordinate several specialized agents---together with tools, memory, and orchestration logic---to solve tasks beyond the reach of any single model, and they now drive applications in coding, scientific discovery, and complex real-world problem solving. This sophistication comes at the cost of reliability: empirical audits report failure rates as high as $86\%$ in popular frameworks, with breakdowns ranging from faulty task decomposition to inter-agent miscommunication~\citep{mast,agentracer}. When such a system fails, the first question a developer must answer is which agent, at which step, committed the \emph{decisive error}---the earliest action whose correction would have averted the failure~\citep{whoandwhen}. Answering it is the prerequisite for any targeted fix, and doing so by hand means combing through long, intertwined interaction logs, an effort that scales poorly as systems grow.

% problem introduction
Multi-agent systems (MAS) powered by large language models (LLMs) have emerged as a transformative paradigm across a wide range of domains, including coding~\citep{qian2023chatdev,hong2023metagpt}, scientific discovery~\citep{lu2024aiscientist,ghafarollahi2024sciagents}, and general-purpose real-world problem solving~\citep{wu2024autogen,fourney2024magenticone}. Yet as these systems grow in both capability and complexity, ensuring their reliability in deployment remains a critical open challenge~\citep{pan2025map,mast}. In particular, when a system fails to complete a task, pinpointing the precise step and responsible agent that first triggered the failure within a long execution trajectory, a problem known as \emph{failure attribution}~\citep{whoandwhen}, is a crucial first step toward monitoring, diagnosing, and ultimately improving these systems~\citep{epperson2025interactive}.

Unfortunately, automated failure attribution in MAS remains an open challenge. Existing studies broadly follow two strategies. LLM-as-a-judge approaches locate the decisive step in under 10\% of cases in the hardest setting~\citep{whoandwhen}, and later adjacent work~\citep{echo, chief, falat, correct} leaves accuracy suboptimal while depending on proprietary models and many inference calls. Training a dedicated attributor via reinforcement learning~\citep{agentracer, graphtracer} improves on prompting but demands compute-intensive training and suffers from data scarcity. Beyond these limitations, both strategies share two deeper problems. First, they treat a trajectory purely as text, leaving the rich signal in model internals unused despite hidden representations being freely inspectable in open-weight models. Second, no existing approach models how the failure signal propagates: since each step conditions on prior outputs, a single upstream error corrupts all subsequent steps~\citep{gao2025singleormulti}, making downstream steps appear erroneous merely by inheritance. Distinguishing root-cause from inherited failure thus requires explicitly modeling this downstream contamination, which is difficult from text alone and motivates looking into model internals. Driven by these two challenges, we introduce \textbf{Causal Residual Rescoring (\crr{})}, a training-free framework that operates on model internals rather than surface text and proceeds in two phases: local scoring, then rescoring with a discounting mechanism.

The local scoring phase draws the signal directly from hidden states. For each step, we compute a per-step anomaly score from the layer activations of an off-the-shelf proxy model, one that did not generate the trajectory, requiring no labels and no access to the generating system's internals. This rests on our observation that decisive error steps often coincide with internally atypical generations, echoing findings in reliable machine learning that internal signals such as activations can capture abnormality~\citep{gradnorm, sal}. Concretely, we factorize the step representations using singular value decomposition (SVD), which helps identify latent subspaces that capture the structure shared by typical generations. We then surface error steps by measuring the departure of each step's representation from this subspace.

The rescoring phase then addresses the downstream contamination problem of the local scores, as a step conditioned on a corrupted context may appear highly anomalous without being the decisive error. \crr{} disentangles the decisive errors from inherited ones by subtracting from each step's score the portion inherited from causally upstream steps. The key insight is that the proxy's own attention reveals how much each step borrows from its predecessors: converting the attention mass a step routes back to earlier steps into dependency weights yields a per-step estimate of inherited anomaly, subtracted under a discount strength. After rescoring, intrinsic anomaly dominates at its origin while the contaminated tail is reduced, exposing the actual root-cause step. 
% \crr{} introduces one continuous hyperparameter, needs no training data or per-trajectory calibration, and wraps any per-step base scorer without modification.

Extensive experimental results on contemporary LLMs confirm that \crr{} effectively improves performance across diverse datasets spanning different multi-agent systems (Section~\ref{sec:exp:main}). Compared to state-of-the-art methods, we consistently improve decisive-error detection accuracy at both the agent and step level on the challenging Who\&When~\citep{whoandwhen}, TraceElephant~\citep{traceelephant}, and CORRECT-Error~\citep{correct} benchmarks. Specifically, on the representative benchmark Who\&When~\citep{whoandwhen}, \crr{} improves step-level accuracy by up to X\% over the second-best method and achieves state-of-the-art performance of Y\%. Furthermore, we conduct comprehensive ablations on the key components of our methodology (Section~\ref{sec:exp:analysis}) and extend our inquiry to showcase \crr{}'s generalization across different model scales and cross-dataset settings (Section~\ref{sec:exp:main}). To summarize our key contributions:

\begin{itemize}
  \item We identify and quantify \emph{downstream error contamination}: because a decisive error corrupts every subsequent step, a local scorer is biased toward steps \emph{after} the true error.
  \item We propose \crr{}, a training-free and generator-agnostic framework that recasts failure attribution as \emph{unsupervised per-step anomaly scoring} over proxy-model representations. \crr{} pairs a local anomaly scorer with an attention-weighted causal rescoring layer that subtracts each step's inherited anomaly to recover the earliest decisive step. Despite using no labels, it is competitive with methods that require fine-tuning or strong LLM judges (\S\ref{sec:method}).
  \item Extensive experiments across multiple proxy models and three failure-attribution benchmarks (Who\&When, TraceElephant, and CORRECT-Error) demonstrate that \crr{} consistently improves both agent- and step-level attribution at no additional supervision cost, and analysis confirms that the gains generalize across model scales and cross-dataset settings (\S\ref{sec:experiments}).
\end{itemize}

% We, therefore, introduce local scoring each step by an anomaly signal read directly from the hidden states of an off-the-shelf proxy model, one that did not generate the trajectory, requiring no labels, no fine-tuning, and no access to the generating system's internals. Casting attribution as a score field also brings a previously unnamed obstacle into view. A decisive error does not stay local: it corrupts the context that every later step is built on, so any scorer that conditions on context---an LLM judge as much as our representation score---registers elevated anomaly throughout the downstream suffix, not only at its source. The naive predictor that takes the argmax therefore drifts past the true error onto a step that is merely contaminated by it.


% We propose Causal Residual Rescoring (\crr{}), which removes from each step's score the portion of its anomaly inherited from causally upstream steps. We estimate inheritance from the proxy's own attention: the attention mass a step routes back into each predecessor, converted into a row-stochastic dependency weight, indicates how much of that step's state was borrowed from earlier ones. Subtracting this attention-weighted upstream anomaly, modulated by a single discount strength, leaves the intrinsic anomaly to dominate at its origin while collapsing the contaminated tail. \crr{} adds one continuous hyperparameter, needs no training data and no per-trajectory calibration, and wraps any per-step base scorer without modification.


% A natural alternative is to score every step rather than guess one---assign each step a scalar reflecting how error-like it is, and read off the decisive step as the maximum. Yet this per-step view has gone largely unexplored, and the reason is instructive. The obvious way to obtain such scores would be a process reward model trained on per-step correctness, as used for mathematical reasoning~\citep{mathshepherd,processbench}; but the supervision available for failure attribution is a single decisive-step label per trajectory---a localization target, not the dense process labels a reward model needs. The data that exists feeds the full-trajectory-to-one-step paradigm, and no per-step process scorer has been built for this task. We sidestep supervision altogether. Borrowing from representation-based out-of-distribution detection~\citep{gradnorm,sal}, we score each step by an anomaly signal read directly from the hidden states of an off-the-shelf proxy model---one that did not generate the trajectory---requiring no labels, no fine-tuning, and no access to the generating system's internals. This unsupervised per-step scorer is no strawman: on its own it already rivals the training-based and advanced-prompting methods above, and as a byproduct it returns a ranked shortlist of suspect steps rather than a lone guess. \TODO{Change the narrative to the internal signal perspective}. 



% However, as MAS increasingly scale in both complexity (e.g., sophisticated interactions with other agents and tools) and deployments, decisive error recognition---pinpointing the precise agent and step that first triggered the failure~\citep{whoandwhen}---has remained a fundamental challenge~\citep{pan2025map}. Unlike SAS where failures can be traced to a single faulty output, MAS failures often emerge from cascading effects across agents: an error-prone step by one agent can propagate through downstream interactions, ultimately causing task failure~\citep{gao2025singleormulti}. Efficient decisive error recognition is essential for reliable MAS deployments, sustaining service quality, and guiding operational management such as safe agent upgrades, targeted restarts, and service monitoring~\citep{epperson2025interactive}.

% Unfortunately, decisive error recognition in MAS remains open-ended due to three fundamental obstacles: (1) \emph{Generality}: MAS span a wide spectrum of applications, and even within an application, requests can exhibit drastically different error patterns, making it difficult to design methods that generalize. Existing advances~\citep{zhang2025agent} often resort to LLM-as-a-judge methods for generality, yet achieve $\leq$10\% accuracy. Recent efforts~\citep{ge2025introducing} to improve accuracy through fine-tuning not only hurt generalization, but fall short in (2) \emph{Data Efficiency}: obtaining labeled data in error recognition is notoriously expensive and inherently ambiguous. For example, in the \textsc{Who\&When} dataset~\citep{zhang2025agent}, annotators spent over 30 expert hours labeling fewer than 200 trajectories, yet disagreement rates exceeded 50\%. This makes any training-based approaches ineffective without voluminous data; and (3) \emph{Computation Efficiency}: even with sufficient data, tuning LLMs for error recognition may be impractical to catch up with deployments where new error types arise continuously and sporadically, e.g., new attacks in cloud AIOps~\citep{intentnet-sigcomm}.

% Automated failure attribution has so far followed two broad strategies, both of which read the full trajectory and return a single point verdict. The first prompts an LLM as a judge: the benchmark study that introduced the task evaluates all-at-once, step-by-step, and binary-search prompting~\citep{whoandwhen}, and later work makes the judge more elaborate---ECHO aggregates several role-conditioned analysts through confidence-weighted voting~\citep{echo}, while CORRECT retrieves distilled error schemata from a cache of past failures to guide the judgment~\citep{correct}. The second strategy fine-tunes a dedicated attributor: AgenTracer trains an 8B tracer with reinforcement learning on thousands of synthetically annotated trajectories~\citep{agentracer}. These methods improve over naive prompting, but their accuracy is purchased with resources that are scarce in this setting. Annotated failure traces are expensive and ambiguous to produce---labeling fewer than two hundred trajectories for Who\&When took over thirty expert hours, with annotator disagreement exceeding fifty percent~\citep{correct}---so schema libraries and fine-tuning alike inherit a heavy labeling burden, while the strongest prompting pipelines depend on proprietary judges and many inference calls. Even then, step-level accuracy stays low: frontier reasoning models locate the decisive step in under ten percent of cases on the harder systems~\citep{whoandwhen,agentracer}.

% A natural alternative is to score every step rather than guess one---assign each step a scalar reflecting how error-like it is, and read off the decisive step as the maximum. Yet this per-step view has gone largely unexplored, and the reason is instructive. The obvious way to obtain such scores would be a process reward model trained on per-step correctness, as used for mathematical reasoning~\citep{mathshepherd,processbench}; but the supervision available for failure attribution is a single decisive-step label per trajectory---a localization target, not the dense process labels a reward model needs. The data that exists feeds the full-trajectory-to-one-step paradigm, and no per-step process scorer has been built for this task. We sidestep supervision altogether. Borrowing from representation-based out-of-distribution detection~\citep{gradnorm,sal}, we score each step by an anomaly signal read directly from the hidden states of an off-the-shelf proxy model---one that did not generate the trajectory---requiring no labels, no fine-tuning, and no access to the generating system's internals. This unsupervised per-step scorer is no strawman: on its own it already rivals the training-based and advanced-prompting methods above, and as a byproduct it returns a ranked shortlist of suspect steps rather than a lone guess. \TODO{Change the narrative to the internal signal perspective}

% Casting attribution as a score field also brings a previously unnamed obstacle into view. A decisive error does not stay local: it corrupts the context that every later step is built on, so any scorer that conditions on context---an LLM judge as much as our representation score---registers elevated anomaly throughout the downstream suffix, not only at its source. The naive predictor that takes the argmax therefore drifts past the true error onto a step that is merely contaminated by it. We argue this is the overlooked reason step-level accuracy lags so far behind agent-level accuracy in prior work: scorers correctly sense that something went wrong but misplace \emph{when}. The pathology is measurable---the offset between predicted and true step is sharply biased toward positive values---and, crucially, it is correctable, because a score field, unlike a single verdict, is something one can recompute. \TODO{Revise, stating the problem more sharply}

% We propose Causal Residual Rescoring (\crr{}), which removes from each step's score the portion of its anomaly inherited from causally upstream steps. We estimate inheritance from the proxy's own attention: the attention mass a step routes back into each predecessor, converted into a row-stochastic dependency weight, indicates how much of that step's state was borrowed from earlier ones. Subtracting this attention-weighted upstream anomaly, modulated by a single discount strength, leaves the intrinsic anomaly to dominate at its origin while collapsing the contaminated tail. \crr{} adds one continuous hyperparameter, needs no training data and no per-trajectory calibration, and wraps any per-step base scorer without modification.

% Our contributions are as follows:
% \begin{itemize}
%   \item We recast failure attribution as \emph{unsupervised per-step anomaly scoring} over proxy-model representations---training-free, annotation-free, and generator-agnostic---and show that this label-free scorer is already competitive with methods that require fine-tuning or strong LLM judges, while additionally yielding a ranked candidate set (\S\ref{sec:method}).
%   \item We identify \emph{downstream error contamination} as a general obstacle that biases any contextual scorer toward steps after the decisive error, explaining the persistent gap between agent- and step-level accuracy, and show it is both measurable and correctable (\S\ref{sec:contam}).
%   \item We propose \crr{}, an attention-weighted causal rescoring layer that subtracts inherited anomaly to recover the earliest decisive step, and demonstrate across base scorers, proxy models, and both Who\&When subsets that it consistently improves step-level attribution at no additional supervision cost (\S\ref{sec:experiments}).
% \end{itemize}